
\subsubsection{2.1 Benchmark Lifecycle and Displacement Dynamics}

The lifecycle of ML benchmarks has received increasing attention as the community grapples with benchmark saturation and the challenge of meaningful progress measurement. Ott et al. \cite{ott2022benchmark} analyze a large collection of benchmarks, finding that benchmark breadth correlates with longevity at the population level. This population-level correlation motivates investigation of whether individual-benchmark breadth predicts displacement timing --- a finer-grained survival analysis question that population-level analyses do not address.

Koch et al. \cite{koch2021reduced} establish a 87-task taxonomy for Papers With Code and measure concentration dynamics in benchmark adoption. Their taxonomy provides the h-e2 panel underlying the present analysis. Where Koch et al. characterize concentration trends across tasks, this paper uses the same panel to test individual predictors of displacement events.

Paullada et al. \cite{paullada2021data} provide a qualitative governance framework for dataset lifecycles, arguing that benchmark datasets are sociotechnical artifacts whose trajectories are shaped by community norms, institutional pressures, and competitive incentives. The present work operationalizes these qualitative observations into a quantitative survival analysis framework.

\subsubsection{2.2 Benchmark Overuse and Community Replacement}

Recent work on benchmarking practices \cite{iclr2025benchmarking} foregrounds benchmark overuse as a driver of community-initiated replacement. Under this framing --- corresponding to the saturation mechanism H2 --- a benchmark that has been widely adopted signals consensus that the task is solved, incentivizing proposals for harder successors. The GLUE $\rightarrow$ SuperGLUE transition is a canonical example: broad community adoption of GLUE \cite{wang2018glue} was followed by recognition of GLUE's ceiling, driving rapid adoption of SuperGLUE \cite{wang2019superglue}.

The null result reported here partially challenges this framing at the individual benchmark level: even if overuse drives replacement in high-profile cases, submission-count breadth at introduction year does not predict displacement timing across the broader population of 87 tasks.

\subsubsection{2.3 Community Adoption and Lock-In Theory}

Rogers' Diffusion of Innovations \cite{rogers2003diffusion} provides the theoretical foundation for the lock-in mechanism H1: early adopters become opinion leaders whose investment in a technology creates network effects that slow transitions to alternatives. The null result here suggests that this investment, at the aggregate level measured by submission count, does not translate into reduced displacement hazard.

\subsubsection{2.4 Survival Analysis in Bibliometrics}

CoxPHFitter from the lifelines library \cite{davidson2019lifelines} is used for all Cox regressions. The L2 penalizer (penalizer = 0.1) was selected for convergence stability. The key methodological challenge --- identifying time-independent predictors that do not collapse into temporal proxies --- is directly addressed by the FAIL FAST protocol.

